\section*{Research programme and collaboration}
This paper grew out of an eight-team project involving 17 students with complementary backgrounds in mathematics, data science, engineering, economics and finance. The size of the author group reflects a division of research responsibilities and a shared learning process: Teams 1--3 developed the mathematical, econometric and simulation foundations; Teams 4--5 built the operating and corporate-valuation layers; and Teams 6--8 studied gold-price dynamics, market evidence and option-implied calibration.

The extended \href{https://raw.githubusercontent.com/StartingFinanceClubPoliTo/Research/main/Barrick-Gold/Barrick\%20Mining\%20UNIFICATO/thesis/Articolo.pdf}{technical report} records definitions, derivations and implementation choices to establish a common technical language across these backgrounds. Teams worked on dedicated modules, then compared assumptions, revised their contributions and reconciled dates, units and interfaces before integration. The present paper condenses this work into a single auditable experiment. Together with the reproducible code, it demonstrates the group's existing research practice and provides a basis for developing a sustained student-led quantitative research team with academic guidance.

\section{Introduction}
A gold producer's value depends on the distribution of future commodity prices as well as its operating and financing assumptions. A single price path can hide this source of model risk. We ask how an unconstrained Barrick operating-proxy distribution changes when progressively richer gold-price processes are substituted within a common operating scenario and discounted-margin contract.

We compare Black--Scholes, Heston, Bates--Poisson and Full Bates--Hawkes calibrated to GLD options. Current-surface diagnostics and rolling out-of-sample tests distinguish fit from predictive performance; shared production, costs and WACC shocks isolate the effect of the gold engine. The contribution is this controlled comparison and its auditable integration, rather than a claim that any model provides an unconditional fair value.
